\section{Exact finite-sample structure}
\label{sec:finite}

We include the elementary likelihood calculation from the preceding
hidden-form report \cite{hidden}, because it makes the statistical conclusion
self-contained and specifies exactly which part of the proof is new.

\begin{lemma}[Counts and likelihood]
\label{lem:likelihood}
Let $M_n$ be the number of nondegenerate projective alternating forms on
$\F^n$. Put
\[
 c_n=\prod_{j=2}^{n/2}(1-q^{-(2j-1)}),\qquad
 p_n=\frac{q^{n-2}-1}{q^{n-1}-1}.
\]
Then
\begin{equation}
 M_n=q^{D-1}c_n,\qquad
 \frac{dP_{n,t}}{dQ_{n,t}}=\frac{Z}{M_np_n^t},\qquad
 \E_Q Z=M_np_n^t.
 \label{eq:finite-likelihood}
\end{equation}
The posterior distribution of $[H]$ under the planted model is uniform on
the $Z$ compatible nondegenerate projective forms. Consequently the exact
Bayes optimal recovery probability is
\begin{equation}
 \operatorname{Rec}_{n,t}=\frac{Q_{n,t}(Z>0)}{M_np_n^t}.
 \label{eq:finite-recovery}
\end{equation}
\end{lemma}

\begin{proof}
The general linear group acts transitively on nondegenerate alternating
forms, with stabilizer $\Sp_n(q)$. Counting ordered bases and symplectic
bases gives
\[
 |\GL_n(q)|=q^{n^2}\prod_{j=1}^n(1-q^{-j}),\qquad
 |\Sp_n(q)|=q^{n(n+1)/2}\prod_{j=1}^{n/2}(1-q^{-2j}).
\]
Their quotient is
$q^D\prod_{j=1}^{n/2}(1-q^{-(2j-1)})$. Dividing by $q-1$ proves the
formula for $M_n$.

For each nonzero $u$, the form $H(u,\cdot)$ is a nonzero linear functional.
Among the $q^n-q$ vectors $v$ independent of $u$, precisely $q^{n-1}-q$
satisfy $H(u,v)=0$. Therefore a uniform two-plane is $H$-isotropic with
probability $p_n$. Every compatible $[H]$ gives likelihood $p_n^{-t}$
relative to $Q$, and incompatible forms give likelihood zero. Averaging over
the $M_n$ forms proves \eqref{eq:finite-likelihood} and the uniform posterior.
Conditional optimal recovery has probability $1/Z$ under $P$, so
\[
 \E_P\frac1Z
 =\E_Q\left[\frac{Z}{M_np_n^t}\frac{\ind_{\{Z>0\}}}{Z}\right]
 =\frac{Q(Z>0)}{M_np_n^t},
\]
where $\ind_{\{Z>0\}}/Z$ is defined to be zero when $Z=0$.
All the basis counts and orthogonality statements are valid in characteristic
two; no division by two is involved.
\end{proof}

The full posterior depends on the observations only through $N$. The
likelihood for testing $P$ against $Q$ depends only on $Z$. The identities of
the candidate forms are still needed to output a recovered form; a bare
count does not encode those identities.

At $t=d+a$,
\begin{equation}
 M_np_n^t=c_nq^{-a}(qp_n)^t\longrightarrow c_qq^{-a}=\lambda_a,
 \label{eq:mean-limit}
\end{equation}
because $qp_n=1-O_q(q^{-n})$ and $t=O(n^2)$. Notice also that
\begin{equation}
 M_np_n^{d+a}<c_nq^{-a}<1\qquad(a\ge0).
 \label{eq:threshold-below-one}
\end{equation}
The likelihood ratio therefore exceeds one exactly when $Z>0$ for every
finite $n\ge4$ at these offsets. What remains to be determined is the
probability of that event and the full law of $Z$.
